\documentclass[a4paper,12pt]{article}
\usepackage[utf8]{inputenc}
\usepackage[T1]{fontenc}
\usepackage[brazil]{babel}
\usepackage{geometry}
\usepackage{setspace}
\usepackage[breaklinks=true, hypertexnames=false]{hyperref}
\usepackage{xurl}
\usepackage{graphicx}
\usepackage{titlesec}
\usepackage{enumitem}
\usepackage{float}
\usepackage{amsmath}
\usepackage{booktabs}
\geometry{margin=2.5cm}
\setstretch{1.3}
\titleformat{\section}{\normalfont\Large\bfseries}{\thesection}{1em}{}
\titleformat{\subsection}{\normalfont\large\bfseries}{\thesubsection}{1em}{}

\title{Estruturação do Sistema MyBook --- E-commerce de Livros}
\author{Iago Lima}
\date{Julho de 2026}

\begin{document}

\begin{titlepage}
    \centering
    % SUBSTITUIR pelo arquivo real da logo, se disponível
    \includegraphics[width=0.35\textwidth]{logo_iff.png}\\[0.5cm]
    {\LARGE \textbf{Instituto Federal de Educação, Ciência e Tecnologia Fluminense} \\
    \vspace{0.5cm}
    \Large Campus Bom Jesus do Itabapoana}\\[3cm]
    {\huge \textbf{Estruturação do Sistema MyBook\\[0.3cm] E-commerce de Livros}}\\[1.5cm]
    {\large Relatório Técnico --- Fase 1: Estruturação do Sistema}\\[4.5cm]
    {\large \textbf{Aluno:} Iago Lima \\[0.3cm]
    \textbf{Disciplina:} Projeto de Sistemas de Software \\[0.3cm]
    \textbf{Professor:} [NOME DO PROFESSOR]\\[0.3cm]
    \textbf{Data:} Julho de 2026}\\[2cm]
    \vfill
\end{titlepage}

\newpage
\tableofcontents
\newpage

\section{Introdução}

O presente relatório documenta a primeira fase do desenvolvimento do sistema
\textbf{MyBook}, um e-commerce de livros voltado ao mercado brasileiro. Nesta fase,
denominada \emph{estruturação do sistema}, foram tomadas as decisões arquiteturais
fundamentais que orientam todo o restante do projeto: a seleção das tecnologias de
implementação, a definição do modelo de hospedagem, o posicionamento quanto ao uso de
Inteligência Artificial (IA) e a caracterização do ambiente de desenvolvimento.

Um e-commerce é uma aplicação de missão crítica que manipula dados sensíveis de
clientes (identificação, endereços e informações de pagamento) e realiza transações
financeiras. Por essa razão, as escolhas estruturais não podem ser arbitrárias:
devem privilegiar segurança, reprodutibilidade do ambiente, portabilidade entre
máquinas e clareza de manutenção. A estruturação adequada nesta fase reduz o custo de
correções futuras e estabelece as fronteiras de responsabilidade de cada componente.

O sistema MyBook adota uma arquitetura \emph{monolítica modular} com renderização no
lado do servidor (\emph{Server-Side Rendering}, SSR). A aplicação é organizada em
camadas de responsabilidade única --- repositórios de acesso a dados, serviços de
regra de negócio, rotas de controle e \emph{views} de apresentação --- o que favorece
a coesão interna e o baixo acoplamento entre módulos. Todo o ambiente é
\emph{containerizado} com Docker, garantindo que a aplicação se comporte de forma
idêntica em qualquer máquina, do computador de desenvolvimento ao servidor de produção.

Este documento está organizado de modo a responder objetivamente às quatro questões
norteadoras da fase de estruturação, apresentando não apenas as decisões tomadas, mas
também a justificativa técnica de cada uma delas.

\section{Objetivos}

\begin{itemize}[leftmargin=2em]
    \item Definir o conjunto de tecnologias (linguagens, frameworks e bibliotecas)
    utilizadas na implementação do sistema.
    \item Especificar o modelo de hospedagem da aplicação, contemplando ambiente de
    desenvolvimento e de produção.
    \item Determinar o posicionamento do projeto quanto ao uso de Inteligência
    Artificial, definindo se seria empregado modelo local ou API externa.
    \item Caracterizar o ambiente de desenvolvimento, incluindo IDE, contêineres de
    isolamento, banco de dados e ferramentas auxiliares.
\end{itemize}

\section{Metodologia}

\subsection{Materiais Utilizados}

Foram utilizados, na estruturação e implementação do sistema, os seguintes recursos de
software e infraestrutura:

\begin{itemize}[leftmargin=2em]
    \item \textbf{Node.js} (versão $\geq$ 20 LTS) --- ambiente de execução JavaScript
    no servidor.
    \item \textbf{Express} 4.x --- framework web para roteamento e \emph{middlewares}.
    \item \textbf{EJS} 6.x --- \emph{template engine} para renderização de páginas no
    servidor (SSR).
    \item \textbf{PostgreSQL} 16 (imagem \texttt{postgres:16-alpine}) --- sistema
    gerenciador de banco de dados relacional.
    \item \textbf{Docker} e \textbf{Docker Compose} v2 --- containerização e
    orquestração dos serviços.
    \item \textbf{Visual Studio Code} com a extensão \emph{Dev Containers} --- ambiente
    de desenvolvimento integrado.
    \item Bibliotecas de apoio: \texttt{pg} (\emph{driver} PostgreSQL),
    \texttt{bcryptjs} (\emph{hash} de senhas), \texttt{express-session} e
    \texttt{connect-pg-simple} (sessões persistidas), \texttt{helmet} (cabeçalhos de
    segurança), \texttt{express-rate-limit} (limitação de requisições) e
    \texttt{multer} (\emph{upload} de arquivos).
    \item Ferramentas auxiliares: Git (controle de versão), GitHub CLI, e cliente de
    banco DBeaver/pgAdmin para inspeção do PostgreSQL.
\end{itemize}

\subsection{Procedimento Experimental}

A estruturação do sistema seguiu as etapas descritas a seguir:

\begin{enumerate}[leftmargin=2em]
    \item Levantou-se o requisito central do projeto: um e-commerce de livros para o
    contexto brasileiro, com catálogo, carrinho, \emph{checkout}, cálculo de frete por
    CEP e pagamento em modo \emph{sandbox}.
    \item Selecionou-se o conjunto de tecnologias com base em maturidade, curva de
    aprendizado e adequação a uma aplicação SSR com banco relacional.
    \item Definiu-se a arquitetura em camadas (repositórios, serviços, rotas,
    \emph{middlewares} e \emph{views}), organizando o código por domínio.
    \item Elaborou-se o \texttt{Dockerfile} multiestágio e o arquivo
    \texttt{docker-compose.yml}, containerizando a aplicação e o banco de dados.
    \item Configurou-se um \emph{dev container} isolado para o ambiente de
    desenvolvimento, com PostgreSQL integrado e ferramentas pré-instaladas.
    \item Estabeleceu-se a política de segurança e de gestão de segredos, centralizando
    todas as credenciais em variáveis de ambiente (\texttt{.env}).
    \item Avaliou-se a necessidade de recursos de Inteligência Artificial no produto,
    concluindo-se pela sua não incorporação em tempo de execução.
\end{enumerate}

\section{Resultados e Discussão}

\subsection{Tecnologias Utilizadas}

O sistema foi implementado sobre a plataforma \textbf{Node.js}, com o framework
\textbf{Express} para a camada web e \textbf{PostgreSQL} como banco de dados
relacional. A escolha por renderização no servidor com \textbf{EJS} dispensa a
complexidade de um \emph{framework} de \emph{frontend} dedicado, entregando páginas
HTML já prontas ao navegador --- decisão coerente com o princípio KISS (\emph{Keep It
Simple}) e favorável ao SEO. A Tabela~\ref{tab:stack} sintetiza a pilha tecnológica
adotada.

\begin{table}[H]
    \centering
    \caption{Pilha tecnológica do sistema MyBook.}
    \label{tab:stack}
    \begin{tabular}{lll}
        \toprule
        \textbf{Camada} & \textbf{Tecnologia} & \textbf{Função} \\
        \midrule
        Runtime          & Node.js $\geq$ 20      & Execução de JavaScript no servidor \\
        Framework web    & Express 4.x            & Roteamento e \emph{middlewares} \\
        Apresentação     & EJS 6.x (SSR)          & Renderização de páginas no servidor \\
        Banco de dados   & PostgreSQL 16          & Persistência relacional \\
        Acesso a dados   & \texttt{pg} 8.x        & \emph{Driver} de conexão \\
        Autenticação     & \texttt{bcryptjs}      & \emph{Hash} de senhas \\
        Sessão           & express-session + PG   & Sessões persistidas no banco \\
        Segurança        & helmet, rate-limit     & Cabeçalhos e limitação de requisições \\
        \emph{Upload}    & multer                 & Recebimento de mídias \\
        Containerização  & Docker + Compose       & Empacotamento e orquestração \\
        \bottomrule
    \end{tabular}
\end{table}

Adicionalmente, dados sensíveis como o CPF do cliente são protegidos em repouso por
criptografia simétrica \textbf{AES-256-GCM}, e todos os formulários que alteram estado
são protegidos contra CSRF (\emph{Cross-Site Request Forgery}). As senhas nunca são
armazenadas em texto claro, sendo submetidas à função de \emph{hash} \texttt{bcrypt}.

\subsection{Hospedagem do Sistema}

O sistema foi projetado para ser \textbf{totalmente containerizado}, o que torna a
hospedagem independente do provedor. A aplicação e o banco de dados são descritos como
serviços no \texttt{docker-compose.yml} e podem ser executados em qualquer ambiente que
disponha de Docker: máquina local, servidor virtual privado (VPS) ou provedor de nuvem.

Em \textbf{desenvolvimento}, o sistema roda localmente em \texttt{http://localhost:3000}.
Em \textbf{produção}, a aplicação é executada por um usuário não-\emph{root} dentro do
contêiner, posicionada atrás de um \emph{proxy} reverso ou balanceador de carga
responsável pela terminação TLS (HTTPS), habilitada pela variável
\texttt{FORCE\_HTTPS=true}. A Tabela~\ref{tab:hosting} resume a configuração de
hospedagem por ambiente.

\begin{table}[H]
    \centering
    \caption{Configuração de hospedagem por ambiente.}
    \label{tab:hosting}
    \begin{tabular}{lll}
        \toprule
        \textbf{Aspecto} & \textbf{Desenvolvimento} & \textbf{Produção} \\
        \midrule
        Execução       & Docker Compose local        & Docker (VPS/nuvem) \\
        Endereço       & \texttt{localhost:3000}     & Domínio atrás de \emph{proxy} TLS \\
        Protocolo      & HTTP                        & HTTPS (\texttt{FORCE\_HTTPS=true}) \\
        Banco          & \texttt{postgres:16-alpine} & PostgreSQL 16 gerenciado/contêiner \\
        Usuário        & \texttt{node} (não-root)    & \texttt{node} (não-root) \\
        \emph{Healthcheck} & Endpoint \texttt{/health} & Endpoint \texttt{/health} \\
        \bottomrule
    \end{tabular}
\end{table}

O contêiner da aplicação possui \emph{healthcheck} próprio, que verifica
periodicamente o endpoint \texttt{/health}, permitindo que o orquestrador reinicie o
serviço automaticamente em caso de falha. As credenciais de banco e o segredo de
sessão são injetados exclusivamente por variáveis de ambiente, nunca ficando
codificadas no repositório.

\subsection{Inteligência Artificial: Modelo Local ou API Externa}

Após avaliação, decidiu-se que o \textbf{produto MyBook não incorpora Inteligência
Artificial em tempo de execução}. Todas as funcionalidades sensíveis são resolvidas por
lógica determinística e verificável, o que aumenta a previsibilidade e a auditabilidade
do sistema. Em particular:

\begin{itemize}[leftmargin=2em]
    \item O pagamento opera em modo \emph{sandbox/mock}, com os pontos de integração
    real isolados para substituição futura por um \emph{gateway} de pagamento.
    \item As capas de livros sem imagem cadastrada são geradas como \textbf{SVG no
    servidor}, de forma offline e sem qualquer dependência de serviços de geração de
    imagem por IA.
\end{itemize}

O único uso de IA no projeto ocorre no \textbf{fluxo de desenvolvimento}, e não no
produto: adota-se o assistente \textbf{Claude Code} (Anthropic), que opera via
\textbf{API externa}. Portanto, quanto à questão norteadora, registra-se que
\emph{não há modelo de IA local}; a única IA envolvida é uma \textbf{API externa}
empregada como ferramenta de apoio à programação, executada dentro de um contêiner
isolado (ver Seção~\ref{sec:devenv}), sem acesso a dados de produção. Essa separação
garante que a ausência de IA no produto seja uma decisão de arquitetura, e não uma
limitação circunstancial.

\subsection{Ambiente de Desenvolvimento}
\label{sec:devenv}

O ambiente de desenvolvimento foi padronizado por meio de um \textbf{Dev Container},
recurso do Visual Studio Code que descreve, em código, toda a configuração da máquina de
trabalho. Isso assegura que qualquer desenvolvedor obtenha um ambiente idêntico,
eliminando o clássico problema do ``na minha máquina funciona''. A
Tabela~\ref{tab:devenv} detalha os componentes do ambiente.

\begin{table}[H]
    \centering
    \caption{Componentes do ambiente de desenvolvimento.}
    \label{tab:devenv}
    \begin{tabular}{ll}
        \toprule
        \textbf{Componente} & \textbf{Descrição} \\
        \midrule
        IDE                & Visual Studio Code + extensão \emph{Dev Containers} \\
        Contêiner base     & \texttt{node:20-bookworm} \\
        Isolamento         & \emph{Dev container} (acesso restrito a \texttt{/workspace}) \\
        Banco integrado    & PostgreSQL 16 (\emph{host} \texttt{db}, porta 5432) \\
        Controle de versão & Git + GitHub CLI \\
        Assistente de IA   & Claude Code CLI (API externa, uso isolado) \\
        Cliente de banco   & DBeaver/pgAdmin em \texttt{localhost:5432} \\
        Sistema \emph{host}& Windows 10 \\
        \bottomrule
    \end{tabular}
\end{table}

O \emph{dev container} sobe automaticamente um \textbf{PostgreSQL 16} associado, de modo
que a aplicação se conecta ao banco sem configuração adicional --- as variáveis
\texttt{DB\_HOST}, \texttt{DB\_USER}, \texttt{DB\_PASSWORD} e \texttt{DB\_NAME} já chegam
prontas ao ambiente. Os \emph{scripts} SQL de \texttt{db/init/} (esquema e \emph{seed})
são executados na primeira criação do banco, deixando-o pronto para uso.

O isolamento proporcionado pelo contêiner é o que torna seguro executar o assistente de
IA em modo de permissões liberadas (\emph{bypass}): como o Claude só enxerga a pasta
\texttt{/workspace} e os dados residem em volumes Docker dedicados, não há risco ao
sistema de arquivos do \emph{host}. As senhas empregadas no ambiente de desenvolvimento
são \emph{defaults} explicitamente marcados como não reutilizáveis em produção.

\subsection{Discussão Geral}

As decisões estruturais apresentadas convergem para três princípios: \textbf{segurança
por padrão} (segredos em variáveis de ambiente, \emph{hash} de senhas, criptografia de
dados sensíveis, proteção CSRF e limitação de requisições), \textbf{reprodutibilidade}
(containerização integral e ambiente de desenvolvimento descrito em código) e
\textbf{simplicidade} (arquitetura monolítica modular com SSR, evitando complexidade
prematura). A opção por não embarcar IA no produto reforça a auditabilidade das regras de
negócio, mantendo o comportamento do sistema determinístico e testável.

\section{Conclusão}

A fase de estruturação do sistema MyBook cumpriu integralmente os objetivos propostos.
Definiu-se a pilha tecnológica --- \textbf{Node.js, Express, EJS e PostgreSQL 16} ---,
sustentada por uma arquitetura monolítica modular em camadas e reforçada por bibliotecas
consolidadas de segurança e persistência.

Quanto à hospedagem, o sistema foi projetado de forma \textbf{containerizada com Docker},
tornando-o portável entre a máquina de desenvolvimento e qualquer servidor de produção,
com HTTPS provido por \emph{proxy} TLS e execução sob usuário não-\emph{root}. No tocante
à Inteligência Artificial, concluiu-se que \textbf{o produto não utiliza IA em tempo de
execução}; a única IA presente é uma \textbf{API externa} (Claude Code) empregada apenas
como ferramenta de apoio ao desenvolvimento, isolada em contêiner --- não havendo,
portanto, modelo local.

Por fim, o ambiente de desenvolvimento foi padronizado por meio de um \emph{Dev
Container} no Visual Studio Code, com PostgreSQL 16 integrado e ferramentas
pré-configuradas, garantindo reprodutibilidade e isolamento. Como limitação, registra-se
que o pagamento permanece em modo \emph{sandbox}, cabendo às fases seguintes a integração
com um \emph{gateway} real e a configuração da infraestrutura de produção (domínio,
certificado TLS e monitoramento). As bases estabelecidas nesta fase, contudo, tornam
essas evoluções incrementais e de baixo risco.

\begin{thebibliography}{99}

\bibitem{nodejs}
OPENJS FOUNDATION. \textbf{Node.js Documentation}. Disponível em:
\url{https://nodejs.org/docs}. Acesso em: jul. 2026.

\bibitem{express}
OPENJS FOUNDATION. \textbf{Express --- Fast, unopinionated, minimalist web framework for
Node.js}. Disponível em: \url{https://expressjs.com}. Acesso em: jul. 2026.

\bibitem{postgres}
THE POSTGRESQL GLOBAL DEVELOPMENT GROUP. \textbf{PostgreSQL 16 Documentation}. Disponível
em: \url{https://www.postgresql.org/docs/16/}. Acesso em: jul. 2026.

\bibitem{docker}
DOCKER INC. \textbf{Docker Documentation}. Disponível em:
\url{https://docs.docker.com}. Acesso em: jul. 2026.

\bibitem{owasp}
OWASP FOUNDATION. \textbf{OWASP Top 10 --- Web Application Security Risks}. Disponível em:
\url{https://owasp.org/www-project-top-ten/}. Acesso em: jul. 2026.

\end{thebibliography}

\end{document}
